% theory/ParabolicDiffusion/GEAR_FVPM_Diffusion_Theory.tex

\documentclass[11pt, a4paper, parskip]{scrartcl}
\usepackage[margin=2.7cm]{geometry}
\usepackage{amsmath}
\usepackage{amssymb}
\usepackage{amsthm}
\usepackage{graphicx}
\usepackage[utf8]{inputenc} % Required for Unicode characters
\usepackage{xcolor}         % Required for coloring text
\usepackage{enumitem}       % For customizing list environments
\usepackage{physics}        % For common physics notations like dv, grad, div, etc.
\usepackage{booktabs}       % For clean tables
\usepackage{seqsplit}
\usepackage[hidelinks]{hyperref}

\usepackage[]{natbib}
\bibliographystyle{apalike}

% \code{} is \texttt{} that may break mid-identifier/mid-path across lines
% (long function names, file paths) instead of overflowing the margin --
% used for long \texttt content in running text. Left as plain \texttt{}
% inside captions/section titles, where \seqsplit's \futurelet-based
% scanning cannot survive LaTeX's full expansion of moving arguments.
\newcommand{\code}[1]{\texttt{\seqsplit{#1}}}

% KOMA-Script disables the classic \rm/\bf/\it/... font-switch commands
% by default; the numbered sections use them freely for math-mode units
% ($h_{\rm phys}$, $c_{\rm fast}$, ...), matching the old article.cls
% setup, so restore the standard LaTeX2e old-font-command behaviour here
% rather than editing every occurrence.
\makeatletter
\DeclareOldFontCommand{\rm}{\normalfont\rmfamily}{\mathrm}
\DeclareOldFontCommand{\sf}{\normalfont\sffamily}{\mathsf}
\DeclareOldFontCommand{\tt}{\normalfont\ttfamily}{\mathtt}
\DeclareOldFontCommand{\bf}{\normalfont\bfseries}{\mathbf}
\DeclareOldFontCommand{\it}{\normalfont\itshape}{\mathit}
\DeclareOldFontCommand{\sl}{\normalfont\slshape}{\@nomath\sl}
\DeclareOldFontCommand{\sc}{\normalfont\scshape}{\@nomath\sc}
\makeatother

\setkomafont{title}{\rmfamily\bfseries\boldmath}
\addtokomafont{section}{\rmfamily\bfseries\boldmath}
\addtokomafont{subsection}{\rmfamily\bfseries\boldmath}
\addtokomafont{subsubsection}{\rmfamily\bfseries\boldmath}

\theoremstyle{plain}
\newtheorem{proposition}{Proposition}
\newtheorem{lemma}{Lemma}
\theoremstyle{definition}
\newtheorem{definition}{Definition}

% ---------------------------------------------------------------------------
% Physics macros. Names are chosen to mirror the GEAR\_FVPM\_DIFFUSION source
% (see the Nomenclature appendix for the symbol <-> code-identifier map).
% ---------------------------------------------------------------------------
\newcommand{\Zfrac}{Z}                                  % metal mass fraction
\newcommand{\rhoZ}{\rho \Zfrac}                          % metal mass density
\newcommand{\driver}{q}                                  % diffusion driver
\newcommand{\Kmat}{\mathbf{K}}                            % diffusion tensor
\newcommand{\Kstar}{\mathbf{K}^{\star}}                   % interface diffusion tensor
\newcommand{\Shear}{\mathbf{S}}                           % shear tensor
\newcommand{\Splus}{\mathbf{S}_{+}}                       % regularised (PSD) shear tensor
\newcommand{\Fdiff}{\mathbf{F}}                           % diffusion flux
\newcommand{\kappaC}{\kappa_{0}}                          % chem\_data->diffusion\_coefficient
\newcommand{\kappaP}{\kappa}                              % per-particle chemistry\_data.kappa
\newcommand{\CCFL}{C_{\mathrm{CFL}}}                      % chem\_data->C\_CFL\_chemistry
\newcommand{\psiHLL}{\psi}                                % hll\_riemann\_solver\_psi
\newcommand{\epsHLL}{\epsilon}                            % hll\_riemann\_solver\_epsilon
\newcommand{\gammak}{\gamma_{k}}                          % kernel\_gamma
\newcommand{\Aij}{A_{ij}}                                 % interface area
\newcommand{\xij}{\mathbf{x}_{ij}}                        % interface position
\newcommand{\nunit}{\mathbf{n}}                           % interface unit normal
\newcommand{\dx}{\Delta x}                                % GIZMO particle size
\newcommand{\dt}{\Delta t}                                % explicit timestep

\title{Theoretical Foundations of Parabolic Diffusion in GEAR FVPM}
\subtitle{Theory document: governing equations, FVPM discretisation, and
numerical scheme}
\author{Darwin Roduit}
\date{\today}

\begin{document}

\maketitle

This is the theoretical and numerical reference for the parabolic
(Fickian) variant of the GEAR FVPM metal-diffusion chemistry scheme
(\code{src/chemistry/GEAR\_FVPM\_DIFFUSION/}), configured with
\code{--with-chemistry=GEAR-FVPM-DIFFUSION\_N}. Every equation is
cross-referenced to the C source. The hyperbolic (Cattaneo-relaxation)
companion variant is covered separately in
\code{GEAR\_FVPM\_Hyperbolic\_Diffusion\_Theory.tex} (this directory);
user-facing configuration and parameters are documented in the RTD pages
under \code{doc/RTD/source/SubgridModels/GEAR/diffusion/}. Regenerate
this PDF with \code{./run.sh}.

\tableofcontents

\section{Introduction}

\subsection{Purpose and Scope}
This document is a rigorous theoretical and numerical account of the
\emph{parabolic} metal-diffusion solver implemented in
\code{src/chemistry/GEAR\_FVPM\_DIFFUSION} (shared headers) and
\code{src/chemistry/GEAR\_FVPM\_DIFFUSION/parabolic} (mode-specific headers)
of the GEAR Finite Volume Particle Method (FVPM) chemistry module. Every
equation below is transcribed from, and cross-referenced to, the C source
so that the document can be used as a derivation-level companion to the
code rather than a qualitative overview.

The code also contains a second, \emph{hyperbolic} (Cattaneo/Maxwell-type)
formulation of metal diffusion, selected at compile time by the macro
\code{CHEMISTRY\_GEAR\_FVPM\_HYPERBOLIC\_DIFFUSION} and implemented in the
sibling \code{hyperbolic/} subdirectory. That formulation is out of scope
here; it is mentioned only in Section~\ref{sec:hyperbolic-relation} to
clarify that the parabolic equation solved in this document is recovered
as its relaxation-time-to-zero limit. Details about the hyperbolic solver
can be found in the corresponding latex document, \code{./GEAR\_FVPM\_Hyperbolic\_Diffusion\_Theory.tex}

\subsection{Overview of the GEAR FVPM Framework}
GEAR's chemistry module rides on top of a meshless Finite Volume Particle
Method (Lagrangian, GIZMO/MFM-style) hydrodynamics solver. Each particle
$i$ carries an effective volume $V_i$ and a geometric moment matrix
$\mathbf{B}_i$ (SWIFT field \code{p->geometry.matrix\_E}, already inverted)
that is computed once per particle in the hydrodynamics density loop and
consumed, unmodified, by the chemistry module. The diffusion solver
therefore does not own its own spatial discretisation: it reuses the
host MFM geometry to build kernel-weighted least-squares gradients and
generalised interface areas between neighbouring particles (see
Section~\ref{sec:fvpm-geometry}).

\subsection{Metal Transport in Astrophysical Simulations}
Unresolved turbulent mixing redistributes heavy elements produced by
stars on scales below the hydrodynamic resolution. Because Lagrangian
particle methods do not exchange mass between neighbours by default
(each particle retains its metal budget unless an explicit mixing
operator is added), an explicit sub-grid diffusion term is required to
reproduce a realistic chemical enrichment of the interstellar and
intergalactic medium. This is the role of the module documented here.

\section{Physical Foundations}

\subsection{Governing Differential Equations}

\subsubsection{Conservation of Metal Mass}
Let $\rho$ be the gas density and $\Zfrac$ the mass fraction of a given
metal species (code index \code{metal} $\in [0,$
\code{GEAR\_CHEMISTRY\_ELEMENT\_COUNT}$)$). The transported quantity is the
metal mass density $\rhoZ$. In the absence of sources (star-formation
depletion) or sinks (supernova/stellar-wind enrichment, handled by other
modules), conservation of metal mass advected with the flow and diffused
relative to it reads
\begin{equation}
  \label{eq:conservation_rho_Z}
  \pdv{(\rhoZ)}{t} + \div (\rhoZ\, \mathbf{v}) = -\div \Fdiff ,
\end{equation}
where $\mathbf{v}$ is the gas velocity and $\Fdiff$ is the diffusive
metal-mass flux. The advective term $\div(\rhoZ\,\mathbf v)$ is handled
implicitly by the fact that $\rhoZ$ is carried as a passive scalar by the
Lagrangian particles (their mass and position already advect with the
flow); the chemistry module solves only the remaining, purely diffusive
part in the frame co-moving with each particle.

\subsubsection{Parabolic Diffusion Equation}
Isolating the diffusive part of Eq.~\eqref{eq:conservation_rho_Z} gives
the parabolic diffusion equation actually integrated by the code
(\code{chemistry\_end\_force}, \code{chemistry.h}):
\begin{equation}
  \label{eq:parabolic_diffusion_rho_Z}
  \pdv{(\rhoZ)}{t} = -\div \Fdiff .
\end{equation}
This equation is discretised in finite-volume (conservative) form: each
particle's metal mass $m_Z = \rhoZ\, V$ is updated by summing the flux
exchanged with every neighbour across the shared interface (see
Section~\ref{sec:flux-computation}).

\subsubsection{Fick's Law and the Diffusion Driver}
\label{sec:ficks-law}
The diffusive flux obeys a generalised (tensorial) Fick's law,
\begin{equation}
  \label{eq:ficks_law}
  \Fdiff = -\Kmat \cdot \grad \driver ,
\end{equation}
where $\Kmat$ is a symmetric, positive semi-definite $3\times3$ diffusion
tensor (Section~\ref{sec:diffusion-tensor}) and $\driver$ is a
\emph{diffusion driver} whose identity is not fixed: depending on the
diffusion mode (Section~\ref{sec:diffusion-modes}) the code sets
$\driver = \rhoZ$ or $\driver = \Zfrac$. This choice is physically
consequential, not merely notational: it changes which PDE is actually
solved (see the boxed remark at the end of
Section~\ref{sec:diffusion-modes}).

\subsection{Assumptions and Simplifications}

\subsubsection{Continuum Approximation}
Equations~\eqref{eq:conservation_rho_Z}--\eqref{eq:ficks_law} presuppose
that the gas and metal fields are continuous on scales larger than the
particle spacing, i.e.\ that the FVPM discretisation resolves the scales
of interest.

\subsubsection{Operator Splitting from Other Physics}
Diffusion is solved with operator splitting: hydrodynamics, gravity,
cooling, star formation and feedback are integrated by other SWIFT
modules on their own schedules, and the chemistry module only ever sees
the metal mass field at the start and end of its own sub-step
(Section~\ref{sec:time-integration}).

\subsubsection{Isotropy versus Anisotropy}
$\Kmat$ can be isotropic ($\Kmat = \kappaP \mathbf I$) or anisotropic
(full symmetric tensor built from the local velocity shear). Both cases
are implemented; see Section~\ref{sec:diffusion-modes}.

\subsection{Metal Transport Parameters}

\subsubsection{Metal Mass Fraction}
$\Zfrac = m_Z / m$, with $m$ the particle's total (gas) mass. Computed by
\code{chemistry\_get\_metal\_mass\_fraction} (\code{chemistry\_getters.h}).
It is dimensionless and, physically, bounded in $[0,1]$; the code enforces
this bound only at specific points (Section~\ref{sec:unphysical}), since
the diffusion solver internally tolerates small transient excursions
below zero to control noise amplification (Section~\ref{sec:gradients}).

\subsubsection{Diffusion Coefficient and Tensor}
\label{sec:diffusion-tensor}
The single global input parameter \code{chem\_data->diffusion\_coefficient}
(denoted $\kappaC$ below) is re-used with three distinct physical
meanings depending on the diffusion mode -- a deliberate but easily
mis-read overload that is documented explicitly here. The struct comment
at \code{parabolic/chemistry\_struct.h} notes the resulting unit
change: in \code{isotropic\_constant} mode $\kappaC$ is a true diffusivity
(units of length$^2$/time); in the other two modes it plays the role of a
dimensionless Smagorinsky-like constant $C_S$ multiplying a
resolution-dependent turbulent diffusivity.

\subsubsection{Shear Tensor}
\label{sec:shear-tensor}
For the turbulence-coupled modes, $\Kmat$ depends on the local,
trace-free (compression/expansion removed) velocity shear tensor
$\Shear$, built from a kernel-filtered ("Favre-like") velocity field
$\tilde{\mathbf v}$ rather than the raw particle velocity, in the spirit
of a large-eddy-simulation (LES) closure. Its full construction is given
in Section~\ref{sec:diffusion-modes}.

\subsection{Diffusion Modes}
\label{sec:diffusion-modes}

The behaviour of the solver is controlled by
\code{chem\_data->diffusion\_mode}
(\code{enum chemistry\_diffusion\_mode}, \code{chemistry\_properties.h}),
which selects one of three closures. In all three cases the diffusion
coefficient is computed once per active particle per force sub-step by
\code{chemistry\_get\_diffusion\_coefficient}
(\code{chemistry\_getters.h}), using the physical (proper, not
comoving) filtered density $\bar\rho = a^{-3}\,$\code{p->chemistry\_data.filtered.rho}
and physical smoothing length $h_{\rm phys} = a\,h$.

\subsubsection{Isotropic Constant Diffusion}
\begin{equation}
  \Kmat = \kappaC \, \mathbf I , \qquad \driver = \rhoZ .
\end{equation}
This is the simplest closure: $\kappaC$ is a user-specified constant
diffusivity, independent of the local flow. The resulting PDE
(Section~\ref{sec:pde-per-mode}) is a genuine linear heat equation for
the metal density itself,
\begin{equation}
  \Fdiff = -\kappaC \,\grad(\rhoZ) ,
\end{equation}
useful primarily as a verification test against analytic diffusion
solutions (e.g.\ the spreading of a Gaussian or delta-function peak; see
the \code{MetalDiffusionOnePeak}/\code{MetalDiffusionTwoPeaks} test
problems shipped alongside this module).

\subsubsection{Isotropic Smagorinsky Model}
\begin{equation}
  \kappaP = \kappaC \, \gammak^{2} \, h_{\rm phys}^{2} \, \bar\rho \,
            \norm{\Shear}_F , \qquad
  \Kmat = \kappaP \, \mathbf I , \qquad \driver = \Zfrac ,
\end{equation}
where $\norm{\Shear}_F = \sqrt{\Shear:\Shear}$ is the Frobenius norm of
the shear tensor of Section~\ref{sec:shear-def} and $\gammak$ is the
kernel's compact-support radius factor (SWIFT's \code{kernel\_gamma},
$O(1$--$2)$ and kernel-choice dependent, so that
$h_{\rm bar} = \gammak\, h$ is the true compact-support radius of the
kernel of smoothing length $h$). This is a standard Smagorinsky
sub-grid-scale closure: the effective diffusivity scales as
(filter length)$^2\times$(local strain rate), exactly as an eddy
viscosity would in an LES momentum closure. The flux is
\begin{equation}
  \Fdiff = -\kappaP \,\grad \Zfrac .
\end{equation}

\subsubsection{Anisotropic Gradient Model}
\begin{equation}
  \kappaP = \kappaC \, \gammak^{2} \, h_{\rm phys}^{2} \, \bar\rho , \qquad
  \Kmat = \kappaP \, \Splus , \qquad \driver = \Zfrac ,
\end{equation}
i.e.\ the same resolution-dependent prefactor as the Smagorinsky mode,
\emph{without} the $\norm{\Shear}_F$ factor, because directional
information is now retained in full through the tensor $\Splus$ rather
than collapsed into a scalar norm (Rennehan-style~\cite{Rennehan2021} gradient/shear-aligned
diffusion). The flux is the full matrix contraction
\begin{equation}
  \Fdiff = -\Kmat \cdot \grad \Zfrac
         = -\kappaP\, \Splus \cdot \grad \Zfrac .
\end{equation}
Implemented in \code{chemistry\_get\_physical\_parabolic\_flux}
(\code{chemistry\_getters.h}).

\subsubsection{Construction of the Shear Tensor and its Regularisation}
\label{sec:shear-def}
The raw (traceless, Hubble-flow-corrected) shear tensor is built in
\code{chemistry\_get\_physical\_shear\_tensor}
(\code{chemistry\_getters.h}) from the symmetric part of the
gradient of the filtered velocity $\tilde{\mathbf v}$ (comoving
gradients, converted to physical units by $a^{-2}$: one factor of
$a^{-1}$ from $\mathbf v_{\rm peculiar} = \mathbf v_{\rm comoving}/a$ and
one from $\grad_{\rm phys} = a^{-1}\grad_{\rm comoving}$):
\begin{equation}
  \label{eq:shear_raw}
  S_{ij} = \frac{1}{2}
           \left(\pdv{\tilde v_i}{x_j} + \pdv{\tilde v_j}{x_i}\right)_{\rm phys}
           + H\,\delta_{ij}
           - \frac{1}{3}\,\mathrm{tr}(\Shear')\,\delta_{ij} ,
\end{equation}
where $H$ is the Hubble parameter (added once to the diagonal, since the
Hubble flow is isotropic and $\Shear$ is symmetric) and the final term
removes the trace of $\Shear' \equiv \tfrac12(\grad\tilde{\mathbf v} +
\grad\tilde{\mathbf v}^T) + H\mathbf I$ \emph{after} the Hubble term has
been added, so that pure compression/expansion (which is not turbulence)
never contributes to the diffusion coefficient. The filtered velocity
$\tilde{\mathbf v} = \overline{\rho\mathbf v}/\bar\rho$ is a
mass(Favre)-weighted, kernel-smoothed velocity evaluated over the full
compact support $h_{\rm bar} = \gammak h$ (wider than the density-defining
smoothing length $h$), accumulated pairwise in
\code{runner\_iact\_chemistry}/\code{runner\_iact\_nonsym\_chemistry}
(\code{chemistry\_iact.h}) with a harmonic-mean density weighting
and a fixed damping prefactor
\code{GEAR\_FVPM\_DIFFUSION\_FILTERING\_SMOOTHING\_FACTOR}$= 0.8$
(\code{chemistry.h}) applied to the accumulated kernel sum before the
particle's own self-contribution $\bar\rho_i \mathbf v_i$ is added back
in \code{chemistry\_end\_density} (\code{chemistry.h}). Filtering
$\mathbf v$ over the wider kernel support before differentiating
suppresses single-particle noise in the shear estimate, consistent with
the LES philosophy underlying the Smagorinsky and anisotropic-gradient
closures.

$\Shear$ as constructed in Eq.~\eqref{eq:shear_raw} is symmetric but not
sign-definite. For the anisotropic mode this is a problem: $\Kmat$ must
be positive semi-definite (PSD), otherwise the flux $\Fdiff = -\Kmat
\cdot\grad\Zfrac$ can point \emph{up} the metallicity gradient in some
direction and the diffusion operator would sharpen rather than smooth
metal contrasts, an unconditionally unstable configuration for an
explicit parabolic solver. \code{chemistry\_regularize\_shear\_tensor}
(\code{chemistry\_getters.h}) fixes this by discarding the
negative-eigenvalue subspace of $\Shear$, following Balarac et
al.~\cite{Balarac2013}:
\begin{equation}
  \label{eq:S_plus_def}
  \Splus \;=\; \sum_{k=1}^{3} \max(0,\lambda_k)\, \mathbf e_k \otimes \mathbf e_k ,
\end{equation}
where $(\lambda_k, \mathbf e_k)$ are the eigenpairs of the symmetric
matrix $\Shear$ (obtained from a closed-form $3\times3$ symmetric
Jacobi-type diagonalisation, \code{sym\_matrix\_diagonalise\_3x3\_d}).

\begin{proposition}[Positive semi-definiteness of $\Splus$]
\label{prop:psd}
The tensor $\Splus$ defined by Eq.~\eqref{eq:S_plus_def} is symmetric
and positive semi-definite for any symmetric input $\Shear$.
\end{proposition}
\begin{proof}
Symmetry is immediate since $\mathbf e_k\otimes\mathbf e_k$ is symmetric
for every $k$ and the $\lambda_k$ are real (eigenvalues of a symmetric
real matrix). For positive semi-definiteness, take any
$\mathbf u \in \mathbb R^3$ and expand it in the orthonormal eigenbasis
$\{\mathbf e_k\}$ of $\Shear$ (orthonormality follows from the spectral
theorem applied to the symmetric matrix $\Shear$):
\begin{equation}
  \mathbf u^T \Splus\, \mathbf u
  = \sum_{k=1}^3 \max(0,\lambda_k)\, \mathbf u^T (\mathbf e_k \otimes \mathbf e_k) \mathbf u
  = \sum_{k=1}^3 \max(0,\lambda_k)\, (\mathbf e_k \cdot \mathbf u)^2 .
\end{equation}
Each term in the sum is a non-negative scalar $\max(0,\lambda_k)$ times a
square, hence non-negative; the sum of non-negative terms is
non-negative, so $\mathbf u^T \Splus\,\mathbf u \geq 0$ for all
$\mathbf u$, i.e.\ $\Splus \succeq 0$.
\end{proof}

Because $\Kmat = \kappaP \Splus$ with $\kappaP \geq 0$, Proposition
\ref{prop:psd} guarantees $\Kmat \succeq 0$ in the anisotropic mode by
construction, for every particle, at every timestep, independent of the
local flow topology.

\subsubsection{Different Modes Solve Different PDEs}
\label{sec:pde-per-mode}
\begin{quote}
Substituting the driver choice of each mode into
Eq.~\eqref{eq:parabolic_diffusion_rho_Z} via Eq.~\eqref{eq:ficks_law}
gives three genuinely different governing equations, not merely three
parametrisations of the same one:
\begin{align}
  \text{isotropic\_constant:} &\quad
    \pdv{(\rhoZ)}{t} = \div\!\big(\kappaC\, \grad(\rhoZ)\big) ,
    \label{eq:pde_const}\\
  \text{isotropic\_smagorinsky:} &\quad
    \pdv{(\rhoZ)}{t} = \div\!\big(\kappaP(\mathbf x,t)\, \grad \Zfrac\big) ,
    \label{eq:pde_smag}\\
  \text{anisotropic\_gradient:} &\quad
    \pdv{(\rhoZ)}{t} = \div\!\big(\Kmat(\mathbf x,t) \cdot \grad \Zfrac\big) .
    \label{eq:pde_aniso}
\end{align}
Equations~\eqref{eq:pde_smag}--\eqref{eq:pde_aniso} diffuse the
\emph{intensive} quantity $\Zfrac$, with the local density folded into
the (state-dependent) coefficient $\kappaP \propto \bar\rho$; this is the
standard sub-grid metal-mixing closure used in the SPH/MFM literature
(e.g.\ Shen, Wadsley \& Stinson~\citeyear{ShenWadsleyStinson2010}; Escala et al.~\citeyear{Escala2018}), chosen
precisely because diffusing $\rhoZ$ directly (as in
Eq.~\eqref{eq:pde_const}) would let density gradients alone drive
spurious metal currents even in a chemically homogeneous ($\grad Z = 0$)
medium. Equation~\eqref{eq:pde_const} does not have this property and is
intended for controlled test problems rather than production
astrophysical runs.
\end{quote}

\section{FVPM Spatial Discretisation}
\label{sec:fvpm-geometry}

\subsection{Meshless Particle Representation}
The computational domain is discretised by particles, each carrying an
effective volume $V_i$ (\code{p->geometry.volume}) computed by the
underlying MFM hydrodynamics solver from a kernel partition of unity.
Interactions occur between particles found by the same neighbour search
(compact-support kernel with smoothing length $h$) used by the
hydrodynamics solver; the chemistry module performs no independent
neighbour finding.

\subsection{Least-Squares Gradient Estimator}
\label{sec:gradient-estimator}
For every scalar field $Q$ carried by the particles (metal mass
fraction, metal density, velocity components), the module computes a
kernel-weighted, matrix-corrected gradient estimate
(\code{chemistry\_gradients\_collect}, \code{chemistry\_gradients.h}):
\begin{equation}
  \label{eq:gradient_estimator}
  (\grad Q)_i \;=\; \sum_{j} (Q_i - Q_j)\, \tilde{\boldsymbol\psi}_i(\mathbf x_j) ,
  \qquad
  \tilde{\boldsymbol\psi}_i(\mathbf x_j) = W_{ij}\, \mathbf B_i \cdot \mathbf{dx}_{ij} ,
\end{equation}
with $\mathbf{dx}_{ij} \equiv \mathbf x_i - \mathbf x_j$ (the code's
sign convention), $W_{ij} \equiv W(|\mathbf{dx}_{ij}|, h_i)$ the kernel
value, and $\mathbf B_i = \mathbf E_i^{-1}$ the inverse of the particle's
geometric moment matrix (\code{p->geometry.matrix\_E}), computed and
maintained by the host hydrodynamics module using the standard
MFM/GIZMO second-moment normalisation of Lanson \& Vila~\citeyear{LansonVila2008} and
Hopkins~\citeyear{Hopkins2015},
\begin{equation}
  \label{eq:E_matrix_def}
  E_i^{kl} = \sum_j V_j\, W_{ij}\; \mathrm{dx}_{ij}^{k}\, \mathrm{dx}_{ij}^{l} .
\end{equation}

\begin{proposition}[Linear consistency of the gradient estimator]
\label{prop:gradient-consistency}
If a scalar field is exactly linear in a neighbourhood of particle $i$,
$Q(\mathbf x) = Q_i + \mathbf G\cdot(\mathbf x - \mathbf x_i)$ for a
constant vector $\mathbf G$, then the estimator of
Eq.~\eqref{eq:gradient_estimator} reproduces the true gradient exactly:
$(\grad Q)_i = \mathbf G$.
\end{proposition}
\begin{proof}
For the assumed linear field, $Q_i - Q_j = \mathbf G\cdot(\mathbf x_i -
\mathbf x_j) = \mathbf G \cdot \mathbf{dx}_{ij}$. Substituting into
Eq.~\eqref{eq:gradient_estimator},
\begin{equation}
  (\grad Q)_i = \sum_j (\mathbf G\cdot\mathbf{dx}_{ij})\, W_{ij}\, \mathbf B_i \cdot \mathbf{dx}_{ij}
  = \mathbf B_i \left(\sum_j W_{ij}\, \mathbf{dx}_{ij}\otimes\mathbf{dx}_{ij}\right) \mathbf G ,
\end{equation}
using the identity $(\mathbf a\cdot\mathbf u)(\mathbf M\mathbf u) =
\mathbf M(\mathbf u\otimes\mathbf u)\mathbf a$ for a vector $\mathbf a$,
matrix $\mathbf M$ and vector $\mathbf u$, applied with $\mathbf a =
\mathbf G$, $\mathbf u=\mathbf{dx}_{ij}$ componentwise, and factoring
$\mathbf B_i$ (constant over the sum) out to the left. Comparing the sum
in parentheses to Eq.~\eqref{eq:E_matrix_def} shows it equals
$\mathbf E_i$ up to the volume weight $V_j$ folded into $W_{ij}$ by
construction in the code's normalisation (the kernel sum already carries
the $V_j$ weighting through the SPH/MFM partition-of-unity
normalisation used to build \code{matrix\_E}). Hence
$(\grad Q)_i = \mathbf B_i\, \mathbf E_i\, \mathbf G = \mathbf E_i^{-1}\mathbf E_i\, \mathbf G = \mathbf G$,
since $\mathbf B_i \equiv \mathbf E_i^{-1}$ by definition.
\end{proof}

Whenever $\mathbf E_i$ is ill-conditioned (near-singular particle
distribution, flagged by \code{fvpm\_part\_geometry\_well\_behaved(p)}
returning false when the stored condition number exceeds
\code{const\_gizmo\_max\_condition\_number}~$=100$), the matrix inversion is
abandoned and the code falls back to the plain (first-order, kernel
gradient) SPH estimator
\begin{equation}
  \tilde{\boldsymbol\psi}_i(\mathbf x_j) = -\frac{W'_{ij}}{r_{ij}}\, \mathbf{dx}_{ij} ,
\end{equation}
with $W'_{ij}$ the kernel derivative and $r_{ij} = |\mathbf{dx}_{ij}|$
(\code{chemistry\_gradients.h}). After accumulation over all
neighbours, gradients are normalised by $h^{-d}$ ($d$ the spatial
dimension) in \code{chemistry\_gradients\_finalise}
(\code{chemistry\_gradients.h}), matching the kernel's own
normalisation convention.

\subsection{Interface Geometry: Generalised Face Area}
\label{sec:face-area}
Fluxes are exchanged across a notional interface between every
interacting pair $(i,j)$. Following Hopkins (2017, eq.~7), the
generalised interface area vector is
\begin{equation}
  \label{eq:face_area}
  \mathbf A_{ij} = -X_i\, W_{ij}\, h_i^{-d}\, \big(\mathbf B_i \cdot \mathbf{dx}_{ij}\big)
                   -X_j\, W_{ji}\, h_j^{-d}\, \big(\mathbf B_j \cdot \mathbf{dx}_{ij}\big) ,
\end{equation}
with $X_i = V_i$, $X_j = V_j$ in the well-behaved case. If the two
particle volumes differ strongly,
$|V_i-V_j|/\min(V_i,V_j) > 1.5\, d$ (the \code{GIZMO\_VOLUME\_CORRECTION}
branch, \code{chemistry\_iact.h}), both weights are replaced by
the smoothing-length-weighted blend
$X_i = X_j = (V_i h_j + V_j h_i)/(h_i+h_j)$
to avoid numerically unstable cancellations in
Eq.~\eqref{eq:face_area} across sharp density (and hence volume)
contrasts. If either particle's geometry matrix is ill-conditioned, the
code instead falls back to the isotropic SPH face-area estimate
\begin{equation}
  A_{ij}^{\rm SPH} = -\Big(h_i^{-(d+1)} V_i^2 W'_{ij} + h_j^{-(d+1)} V_j^2 W'_{ji}\Big)/r_{ij} ,
  \qquad \mathbf A_{ij} = A_{ij}^{\rm SPH}\, \mathbf{dx}_{ij} .
\end{equation}
The interface normal is $\nunit = \mathbf A_{ij}/\Aij$ with
$\Aij = |\mathbf A_{ij}|$
(\code{chemistry\_iact.h}).

\subsection{Interface Position and Interface-Frame Boost}
\label{sec:interface-position}
The effective interface point, relative to particle $i$, is (Hopkins
2017, eq.~8)
\begin{equation}
  \label{eq:interface_position}
  \xij = x_{\rm fac}\, \mathbf{dx}_{ij} , \qquad
  x_{\rm fac} = -\frac{h_i}{h_i+h_j} ,
\end{equation}
so that the resolved interface sits at
$\mathbf x_i + \xij = (h_j\, \mathbf x_i + h_i\, \mathbf x_j)/(h_i+h_j)$:
a smoothing-length-weighted point between the two particles, pulled
towards the particle with the \emph{smaller} smoothing length (higher
local resolution). The corresponding interface (fluid) velocity used to
boost the hydrodynamic primitive variables into the interface frame
before the Riemann problem is solved is (Hopkins~\citeyear{Hopkins2017diffusion}, eq.~9)
\begin{equation}
  \mathbf v_{ij} = \mathbf v_i + (\mathbf v_i - \mathbf v_j)\, x_{\rm fac}
  \qquad \text{(\code{chemistry\_iact.h})} .
\end{equation}
This position and boost are computed once per pair and reused
identically for the hydrodynamic reconstruction and for the diffusion
Riemann problem, so that both solvers agree on where and in which frame
the interface lives.

\section{Gradient Reconstruction and Slope Limiting}
\label{sec:gradients}

\subsection{Cell-Wise Extremum-Preserving Limiter}
Before any face reconstruction, gradients are scaled down (never up) so
that a linear extrapolation to the farthest interacting neighbour cannot
create a value outside the range spanned by the particle and all its
neighbours (Hopkins~\citeyear{Hopkins2015} formalism,
\code{chemistry\_slope\_limiters\_cell.h}). For a scalar $Q$ with cell value
$Q_i$, neighbour extrema $Q_{\min}, Q_{\max}$ and maximal neighbour
distance $r_{\max}$ (\code{chi->limiter.maxr}), the limiter factor
$\alpha\in[0,1]$ is (\code{chemistry\_slope\_limit\_quantity},
\code{chemistry\_slope\_limiters\_cell.h})
\begin{align}
  g &= r_{\max}\, |\grad Q|_i , \\
  \beta &= \max\!\Big(\beta_{\min},\ \beta_{\max}\min\!\Big(1,\ \tfrac{N_{\rm cond}}{\kappa(\mathbf E_i)}\Big)\Big) ,
  \label{eq:cell_limiter_beta} \\
  \alpha_0 &= \min\!\left(1,\ \beta \, \min\!\left(\frac{Q_{\max}-Q_i}{g},\ \frac{Q_i - Q_{\min}}{g}\right)\right) ,
\end{align}
with $\beta_{\min}=1$, $\beta_{\max}=2$
(\code{GIZMO\_SLOPE\_LIMITER\_BETA\_MIN}/\code{MAX}) and $N_{\rm cond}=100$
(\code{const\_gizmo\_max\_condition\_number}) a reference condition number:
as the particle's geometry matrix approaches singularity
($\kappa(\mathbf E_i)\to N_{\rm cond}$), $\beta\to\beta_{\min}$ and the
limiter becomes maximally conservative; in a well-conditioned
configuration $\beta\to\beta_{\max}=2$, the standard Hopkins (2015)
value.

For quantities that must remain non-negative under extrapolation
(the diffusion drivers $\rhoZ$ and $\Zfrac$, invoked with
\code{pos\_preserve=1}), a second, tighter bound is applied
(\code{chemistry\_slope\_limiters\_cell.h}):
\begin{align}
  f_{\rm shoot} &= \min\!\big(Q_i - Q_{\min} + \tau_{\rm shoot}(Q_{\max}-Q_i),\ Q_{\max}-Q_i\big) , \\
  Q_{\min}^{\rm floor} &= \max\!\left(f_{\rm shoot},\ \tfrac12(Q_i+Q_{\min})\right) , \\
  \alpha_{\rm pos} &= \frac{Q_i - Q_{\min}^{\rm floor}}{g} , \\
  \alpha &= \min(\alpha_0,\ \alpha_{\rm pos}) ,
\end{align}
where $\tau_{\rm shoot} = 0.2$
(\code{GEAR\_FVPM\_DIFF\_CELL\_LIMITER\_SHOOT\_TOLERANGE}) is an allowed
fractional overshoot that keeps the limiter from stalling the diffusion
front entirely in regions with strong local extrema. If
$Q_{\min}=Q_{\max}=Q_i$ (an isolated local extremum, e.g.\ a lone
enriched particle surrounded by pristine gas) the gradient is forced to
zero exactly ($\alpha=0$), since Eq.~\eqref{eq:cell_limiter_beta}--above
would otherwise return a numerically enormous $\alpha$ from a
$0/0$-adjacent limit.

\textbf{Important asymmetry.} The generic cell limiter
\code{chemistry\_slope\_limit\_cell} (called once at the end of the
gradient loop, \code{chemistry\_gradients.h}) is applied to the
velocity gradients but \emph{deliberately not} to the $\Zfrac$/$\rhoZ$
gradients, per the code comment at
\code{chemistry\_slope\_limiters\_cell.h}: cell-limiting the
diffusion driver's gradient would permanently reduce the magnitude of
$\grad\driver$ used everywhere it appears in
Eqs.~\eqref{eq:pde_const}--\eqref{eq:pde_aniso}, i.e.\ it would silently
alter the PDE being solved rather than merely stabilise its
discretisation. Instead, the driver gradients are cell-limited
\emph{locally, on the fly}, immediately before each face reconstruction,
inside \code{chemistry\_gradients\_predict}
(parabolic-mode state reconstruction, see below) and
\code{chemistry\_gradients\_predict\_Z}
(\code{chemistry\_gradients.h}), using the same
$\alpha$-computation above but re-evaluated at every pairwise
interaction. This keeps the corrected gradient reserved only for the
extrapolation actually being performed, without corrupting the
particle's stored gradient field for later use in the flux formula
(Section~\ref{sec:flux-computation}), which needs the \emph{uncorrected}
gradient.

\subsection{Face-Wise (GIZMO) Limiter}
\label{sec:face-limiter}
After the cell-limited gradient is linearly extrapolated to the
interface point $\xij$ (Eq.~\eqref{eq:interface_position}), giving a
provisional reconstructed value $\phi_{\rm mid}^{0} = \phi_i + \grad\phi_i
\cdot \xij$, a second limiting pass constrains the extrapolated pair
$(\phi_{\rm mid,i},\phi_{\rm mid,j})$ using both particle values, not
just particle $i$'s neighbourhood extrema
(\code{chemistry\_slope\_limit\_face\_quantity\_double},
\code{chemistry\_slope\_limiters\_face.h}):
\begin{align}
  \delta_1 &= \tfrac12 |\phi_i - \phi_j| , \qquad \delta_2 = \tfrac14 |\phi_i-\phi_j| , \\
  \bar\phi &= \phi_i + \frac{|\xij|}{r_{ij}}(\phi_j - \phi_i) , \\
  \phi_{+} &=
    \begin{cases}
      \phi_{\max} + \delta_1 & \text{if } \mathrm{sign}(\phi_{\max}+\delta_1) = \mathrm{sign}(\phi_{\max}) \\
      \phi_{\max}/\big(1+\delta_1/|\phi_{\max}|\big) & \text{otherwise}
    \end{cases} \\
  \phi_{-} &=
    \begin{cases}
      \phi_{\min} - \delta_1 & \text{if } \mathrm{sign}(\phi_{\min}-\delta_1) = \mathrm{sign}(\phi_{\min}) \\
      \phi_{\min}/\big(1+\delta_1/|\phi_{\min}|\big) & \text{otherwise}
    \end{cases}
\end{align}
with $\phi_{\min}=\min(\phi_i,\phi_j)$, $\phi_{\max}=\max(\phi_i,\phi_j)$.
The final, limited extrapolated value is
\begin{equation}
  \phi_{\rm mid} =
  \begin{cases}
    \max\!\big(\phi_-,\ \min(\bar\phi+\delta_2,\ \phi_{\rm mid}^{0})\big) & \text{if } \phi_i < \phi_j \\
    \min\!\big(\phi_+,\ \max(\bar\phi-\delta_2,\ \phi_{\rm mid}^{0})\big) & \text{otherwise,}
  \end{cases}
\end{equation}
and the returned, limited increment is $\phi_{\rm mid} - \phi_i$. The
sign-aware construction of $\phi_\pm$ avoids letting the bound cross zero
when $\phi_{\min}$ or $\phi_{\max}$ is itself close to zero, which
matters for a strictly non-negative field like $\Zfrac$ or $\rhoZ$.
Optionally (used for the hyperbolic scheme only, not here) the result can
be clamped to non-negative values by $\phi_{\rm mid}\to\max(0,\phi_{\rm
mid})$. For parabolic diffusion this exact limiter is applied
identically to $\Fdiff$'s driver reconstruction via
\code{chemistry\_slope\_limit\_face}
(\code{parabolic/chemistry\_slope\_limiter\_face.h}), which simply
calls the routine above once for each side of the interface ($i\to j$
and $j\to i$).

\subsection{State Reconstruction at the Interface}
\label{sec:state-reconstruction}
The full per-pair, per-species reconstruction pipeline,
\code{chemistry\_gradients\_predict}
(\code{parabolic/chemistry\_gradients.h}), proceeds as follows for
the metal-density state $U \equiv \rhoZ$:
\begin{enumerate}
  \item Start from the comoving metal densities $U_i, U_j$
        (\code{chemistry\_get\_comoving\_metal\_density}).
  \item Cell-limit the stored $\grad(\rhoZ)$ gradients of $i$ and $j$
        \emph{on the fly}, with the positivity-preserving branch enabled
        (Section~\ref{sec:gradients}).
  \item Linearly extrapolate to the interface point $\xij$
        (Section~\ref{sec:interface-position}):
        $\delta U_i = \grad(\rhoZ)_i \cdot \xij$,
        $\delta U_j = \grad(\rhoZ)_j \cdot \xij^{(j)}$, with
        $\xij^{(j)} = \xij + \mathbf{dx}_{ij}$ the same physical point
        expressed relative to particle $j$.
  \item Apply the face-wise limiter of Section~\ref{sec:face-limiter} to
        $(\delta U_i,\delta U_j)$ jointly, then set
        $U_i \mathrel{+}= \delta U_i$, $U_j \mathrel{+}= \delta U_j$.
  \item Correct unphysical excursions
        (\code{chemistry\_gradients\_correct\_unphysical\_states},
        \code{parabolic/chemistry\_gradients.h}; see
        Section~\ref{sec:unphysical} for the full logic): if the
        extrapolation would push the implied metal mass negative or
        above the gas mass, the offending state is reverted to its
        pre-extrapolation (0th-order) value.
  \item \textbf{Final safety clamp}: irrespective of the above,
        $U_i \gets \max(0, U_i)$ and $U_j \gets \max(0, U_j)$
        unconditionally, since the Riemann solver
        (Section~\ref{sec:flux-computation}) is not robust to residual
        negative densities surviving floating-point round-off in the
        limiter chain.
  \item Convert $U_i, U_j$ from comoving to physical density,
        $U \to a^{-3} U$.
\end{enumerate}
When the diffusion driver is $\Zfrac$ rather than $\rhoZ$ (Smagorinsky
and anisotropic modes), the analogous routine
\code{chemistry\_gradients\_predict\_Z}
(\code{chemistry\_gradients.h}) performs the same steps on the
metal mass fraction gradients, followed by a hard clamp of the
reconstructed metallicity to $[0,1]$
(\code{chemistry\_check\_unphysical\_metallicity},
Section~\ref{sec:unphysical}).

\section{Flux Computation}
\label{sec:flux-computation}

\subsection{Wave-Speed Estimation}
The Riemann solver used here (\code{chemistry\_riemann\_solver\_hopkins2017\_HLL},
\code{parabolic/chemistry\_riemann\_HLL.h}) is explicitly noted in
the source as valid \emph{for parabolic diffusion only} -- it degrades
for the hyperbolic system when the relaxation time is large. Its wave
speeds are not diffusion signal speeds but ordinary \emph{hydrodynamic}
speeds evaluated at the interface: with $u_L, u_R$ the particle
velocities projected onto the interface normal $\nunit$ and
$c_{s,L}, c_{s,R}$ the (physical) hydrodynamic sound speeds,
\begin{equation}
  \lambda_+ = |u_R - u_L| + \max(c_{s,L}, c_{s,R}) , \qquad \lambda_- = -\lambda_+ .
\end{equation}
This choice reuses the same characteristic-speed estimate as the
GIZMO/MFM Euler solver: it guarantees the diffusion Riemann fan is at
least as wide as the hydrodynamic one, so that the diffusive update never
outruns the CFL condition already enforced by the hydro solver on the
same pair. If $\lambda_+-\lambda_- $ falls below a tolerance
$10^{-8}$ (\code{GEAR\_FVPM\_DIFF\_WAVESPEED\_ESTIMATE\_DIFFERENCE\_TOLERANCE}),
both particles are (numerically) at rest relative to the interface with
no local sound speed and the flux is set to zero to avoid a division by
a near-zero denominator.

\subsection{Interface Diffusion Tensor}
The tensor entering the flux is a simple arithmetic mean of the two
particles' physical diffusion tensors,
\begin{equation}
  \Kstar = \tfrac12\big(\Kmat_i + \Kmat_j\big)
  \qquad(\code{chemistry\_riemann\_compute\_K\_star}) .
\end{equation}
If $\norm{\Kstar}_F = 0$ (e.g.\ at the very first timestep, before
$\bar\rho$ has been computed and $\kappaP$ is still zero) the flux is
set to zero identically.

\subsection{The Hopkins~\citeyear{Hopkins2017diffusion} HLL-Type Flux}
\label{sec:hll-flux}
Let $U \equiv \rhoZ$ (the conserved state) and recall $\driver$ is the
diffusion driver, reconstructed to the interface as $\driver_L,
\driver_R$ (Section~\ref{sec:state-reconstruction}), with
$U^\star = \tfrac12(U_L+U_R)$, $\driver^\star = \tfrac12(\driver_L +
\driver_R)$ and $\grad\driver^\star = \tfrac12(\grad\driver_L +
\grad\driver_R)$ (particle-centred, \emph{not} interface-extrapolated,
gradients of the two neighbours).

The flux of each particle's own local, first-order-reconstructed
diffusive flux (Eq.~\eqref{eq:ficks_law}, evaluated with each particle's
own stored gradient) is projected onto the interface normal,
\begin{equation}
  \mathrm{Flux}_{L} = \Fdiff_L\cdot\nunit , \qquad \mathrm{Flux}_R = \Fdiff_R\cdot\nunit .
\end{equation}
The two standard HLL building blocks are then formed,
\begin{equation}
  F_2 = \frac{\lambda_+\,\mathrm{Flux}_L - \lambda_-\,\mathrm{Flux}_R}{\lambda_+-\lambda_-} ,
  \qquad
  F_U = \frac{\lambda_+\lambda_- \,(U_R-U_L)}{\lambda_+-\lambda_-} ,
\end{equation}
i.e.\ $F_2$ is the upwind-blended transport term and $F_U$ is the
standard HLL numerical-dissipation term. Unlike a textbook HLL solver,
which would simply return $F_2 + F_U$, this solver blends the two with a
\emph{minmod} limiter and an \emph{artificial diffusivity} $\alpha$
(Hopkins~\citeyear{Hopkins2017diffusion}, eqs.~7--11) that caps the numerical dissipation at (a
fraction of) the physically prescribed diffusivity:
\begin{equation}
  \label{eq:flux_hll}
  F_{\rm HLL} = \mathrm{minmod}\Big((1+\psiHLL)\, F_2,\ F_2 + \alpha\, F_U\Big) ,
\end{equation}
with the two-argument minmod of Eq.~\eqref{eq:minmod_def} below and
$\psiHLL = $ \code{chem\_data->hll\_riemann\_solver\_psi} a user-tunable
parameter bounding how much extra transport-term weight is admitted
before the antidiffusive correction is clipped.

\subsubsection{The Artificial Diffusivity $\alpha$}
Computed in \code{chemistry\_riemann\_compute\_alpha}
(\code{chemistry\_riemann\_utils.h}):
\begin{align}
  c_{\rm fast}^\star &= \tfrac12(c_{s,L}+c_{s,R}) , \qquad
  v_{\rm HLL} = \tfrac12|u_R-u_L| + c_{\rm fast}^\star , \\
  r &= \frac{v_{\rm HLL}\, |\mathbf{dx}_{ij}|\, |\driver^\star|}{\norm{\Kstar}_F \, |U^\star|} ,
  \qquad
  r_{\rm term} = \frac{0.2+r}{0.2+r+r^2} , \\
  \mathrm{norm\_term} &=
    \frac{\big\| \Kstar \cdot \grad\driver^\star \big\|}{\norm{\Kstar}_F\,\big\|\grad\driver^\star\big\|}
    \quad\big(\to 1 \text{ if } \grad\driver^\star = 0,\ \text{by convention}\big), \\
  \alpha &= \mathrm{norm\_term}\times r_{\rm term} ,
  \qquad \alpha \to 0 \ \text{ if } U^\star = 0 \text{ or } \norm{\Kstar}_F = 0 .
  \label{eq:alpha_def}
\end{align}
The dimensionless ratio $r$ compares the characteristic HLL numerical
flux scale ($v_{\rm HLL}\,\driver^\star/|\mathbf{dx}_{ij}|$, after
clearing the $|\mathbf{dx}_{ij}|$ factor) against the physical diffusive
flux scale ($\norm{\Kstar}_F\,U^\star/|\mathbf{dx}_{ij}|^2$, similarly
normalised): $r_{\rm term}\to 1$ when $r\gg1$ (numerical transport
dominates -- little artificial diffusivity is needed to stabilise it) and
$r_{\rm term}\to 5r$ as $r\to0$ (numerical transport is weak compared to
the physical diffusivity, so the HLL dissipation term is allowed through
more strongly to avoid an under-diffusive, oscillatory update).
\code{norm\_term} additionally attenuates $\alpha$ when $\Kstar$ is
strongly anisotropic relative to the local gradient direction, since
$\|\Kstar\cdot\grad\driver^\star\| \leq \norm{\Kstar}_F
\|\grad\driver^\star\|$ with equality only when $\grad\driver^\star$ lies
along $\Kstar$'s dominant eigendirection.

\begin{definition}[Two-argument minmod]
\label{eq:minmod_def}
\begin{equation}
  \mathrm{minmod}(a,b) =
  \begin{cases}
    \min(a,b) & a>0 \text{ and } b>0 \\
    \max(a,b) & a<0 \text{ and } b<0 \\
    0 & \text{otherwise.}
  \end{cases}
\end{equation}
\end{definition}

\subsection{Direct-Flux Sign-Consistency Guard}
\label{sec:anti-diffusion-guard}
As a final robustness check, the code independently computes a
\emph{direct} (central-difference) estimate of the flux, built purely
from the two particles' own driver values (not their gradients, and not
interface-reconstructed):
\begin{equation}
  q_i = \max(0,\driver_i) , \quad q_j=\max(0,\driver_j) , \quad
  \grad\driver_{\rm dir} = \frac{(q_j-q_i)}{|\mathbf{dx}_{ij,\rm phys}|^2}\,\mathbf{dx}_{ij,\rm phys} ,
\end{equation}
\begin{equation}
  F_{\rm dir} =
  \begin{cases}
    -\big(\Kstar\cdot\grad\driver_{\rm dir}\big)\cdot \dfrac{\mathbf{dx}_{ij,\rm phys}}{|\mathbf{dx}_{ij,\rm phys}|}
      & \text{anisotropic\_gradient mode} \\[1.2em]
    -\bar\kappa\, \grad\driver_{\rm dir}\cdot \dfrac{\mathbf{dx}_{ij,\rm phys}}{|\mathbf{dx}_{ij,\rm phys}|}
      & \text{otherwise, } \bar\kappa=\tfrac12(\kappaP_i+\kappaP_j) .
  \end{cases}
\end{equation}
If this direct estimate disagrees in \emph{sign} with the HLL flux of
Eq.~\eqref{eq:flux_hll} by more than a relative tolerance
$\epsHLL \equiv$ \code{chem\_data->hll\_riemann\_solver\_epsilon},
\begin{equation}
  F_{\rm dir}\cdot F_{\rm HLL} < 0
  \quad\text{and}\quad
  |F_{\rm dir}| > \epsHLL\, |F_{\rm HLL}| ,
\end{equation}
the interface flux is set to zero rather than trusted
(\code{parabolic/chemistry\_riemann\_HLL.h}). This guards against
the minmod-blended HLL construction of Eq.~\eqref{eq:flux_hll}
occasionally returning an up-gradient (anti-diffusive) flux due to
reconstruction error, at the cost of locally stalling diffusion across
that one interface for one sub-step. The final flux entering the update
is $F_{\rm HLL}\times\Aij$ (physical face area,
\code{parabolic/chemistry\_flux.h}).

\section{Flux Limiting for Physical Realizability}
\label{sec:flux-limiter}

Even a Riemann-consistent flux can, once multiplied by a finite timestep
and summed over all of a particle's neighbours, remove more metal mass
than the particle possesses. \code{chemistry\_limit\_metal\_mass\_flux}
(\code{chemistry\_flux.h}) rescales the raw \emph{pairwise} flux
$\Delta m \equiv F_{\rm HLL}\Aij\, \dt_{\min}$ ($\dt_{\min}$ the smaller
of the two particles' active chemistry timesteps) through a sequence of
symmetric (particle order-independent) guards:

\begin{enumerate}
  \item \textbf{Source constraint.} Let $m_{Z}^{\rm src}$ be the current
        metal mass of whichever particle is donating mass (determined
        by $\mathrm{sign}(\Delta m)$, not by particle index -- this is
        what makes the limiter symmetric under $i\leftrightarrow j$). If
        $m_Z^{\rm src}\le 0$, the flux is annihilated: a particle cannot
        export metal mass it does not have.
  \item \textbf{Capacity constraint.} Let $m^{\rm sink}$ be the sink
        particle's total (gas) mass and $m_Z^{\rm sink}$ its current
        metal mass; define the remaining headroom before $\Zfrac=1$,
        $C \equiv m^{\rm sink} - m_Z^{\rm sink}$. If $C\le 0$ the flux is
        annihilated.
  \item \textbf{Noise gate.} If
        $|\Delta m| < \epsilon_{\rm noise}\, m_Z^{\rm src}$ with
        $\epsilon_{\rm noise}=10^{-15}$
        (\code{GEAR\_FVPM\_DIFF\_NOISE\_GATE}), the flux is treated as
        floating-point noise and zeroed, preventing a slow
        accumulation ("ratchet effect") of sub-machine-precision
        transfers.
  \item \textbf{Smooth stability/startup limiter.} An effective mass
        budget is formed,
        \begin{equation}
          m_{\rm eff} = \min\!\Big(\min(m_Z^{\rm src}, C),\ \
          \frac{\max\!\big(0.25\, m_Z^{\rm sink},\ 0.1\, m_Z^{\rm src}\big)}{0.5}\Big) ,
        \end{equation}
        combining the source/capacity bound with a
        \emph{sink-stability} bound $0.25\, m_Z^{\rm sink}$
        (\code{GEAR\_FVPM\_DIFF\_LIMITER\_SINK\_STABILITY}: no single
        pairwise interaction may change the sink's metal mass by more
        than $25\%$, preventing oscillatory "neighbour flooding") and a
        \emph{startup} floor $0.1\, m_Z^{\rm src}$
        (\code{GEAR\_FVPM\_DIFF\_LIMITER\_STARTUP}: without this floor,
        $m_Z^{\rm sink}=0$ pristine gas would have a zero stability
        bound and could never be enriched, stalling the diffusion
        front). The safety scale $0.5$
        (\code{GEAR\_FVPM\_DIFF\_LIMITER\_SAFETY}) leaves headroom below the
        hard constraints. The flux is then attenuated smoothly rather
        than hard-clipped,
        \begin{equation}
          x = \frac{|\Delta m|}{0.5\, m_{\rm eff} + 10^{-40}} , \qquad
          \Delta m \;\longrightarrow\; \frac{\Delta m}{1+x} .
        \end{equation}
        As $|\Delta m|\to\infty$ this rational form asymptotes to a
        bounded transferred mass $\approx 0.5\,m_{\rm eff}$, while for
        $|\Delta m|\ll 0.5\,m_{\rm eff}$ the attenuation factor is
        $\approx 1$ (unmodified flux), giving a $C^0$-continuous
        transition instead of a hard clip that would otherwise reduce
        the formal order of convergence of the scheme near the
        constraint boundary.
\end{enumerate}

\section{Time Integration}
\label{sec:time-integration}

The scheme is explicit and operator-split from hydrodynamics, gravity,
cooling and star formation/feedback. For each active particle pair, the
signed flux is accumulated into both particles' per-timestep flux
buffers with the sign convention that a positive flux (in the direction
of $\nunit$, from $i$ towards $j$) is subtracted from the left state $i$
and added to the right state $j$
(\code{chemistry\_part\_update\_fluxes\_left/right},
\code{parabolic/chemistry\_flux.h}):
\begin{equation}
  \Delta m_{Z,i} \mathrel{-}= F_{\rm HLL}\,\Aij\, \dt_{\min} , \qquad
  \Delta m_{Z,j} \mathrel{+}= F_{\rm HLL}\,\Aij\, \dt_{\min} .
\end{equation}
Because SWIFT is an individual-timestep code, inactive neighbours of an
active particle must still receive their share of the exchanged flux to
preserve exact pairwise mass conservation; this is guaranteed structurally
by \code{runner\_iact\_chemistry\_fluxes\_common}
(\code{chemistry\_iact.h}), which always updates both sides of a
symmetric interaction. At the end of the force loop, each particle
applies its accumulated buffer to the conserved variable
(\code{chemistry\_end\_force}, \code{chemistry.h}):
\begin{equation}
  m_{Z,i}^{n+1} = m_{Z,i}^{n} + \sum_{j\in\mathcal N(i)} \Delta m_{Z,i}^{(j)} ,
\end{equation}
i.e.\ a first-order (forward Euler) update of the conserved metal mass
using the already-limited, already-integrated-over-$\dt_{\min}$ flux sum:
no further sub-cycling or higher-order time integrator is applied to the
diffusion term itself.

\section{Robustness: Unphysical-State Detection and Correction}
\label{sec:unphysical}

All corrections live in \code{chemistry\_unphysical.h} and are invoked
from exactly three call sites, distinguished by an integer
\code{callloc}, each with a different tolerance for how "wrong" a state
is allowed to be before the code intervenes rather than simply
correcting silently
(\code{chemistry\_check\_unphysical\_state},
\code{chemistry\_unphysical.h}):

\begin{itemize}
  \item \textbf{\code{callloc=0}}, after the density loop
        (\code{chemistry\_end\_density}): a negative metal mass here
        indicates corruption \emph{outside} the diffusion solver (e.g.\
        a bug in feedback/enrichment), so the first occurrence for a
        given particle raises a hard \code{error()}.
  \item \textbf{\code{callloc=1}}, during gradient extrapolation
        (Section~\ref{sec:state-reconstruction}): the correction is
        soft. If the extrapolated mass would go negative, the
        extrapolation is discarded and the pre-extrapolation (0th-order,
        particle-centred) value is used instead; the routine returns a
        \code{char corrected} flag consumed by the caller to also reset
        the reconstructed \emph{state} $U$, not just the mass used for
        the check; the flag is compared with an integer test rather than
        a floating-point comparison, avoiding round-off-triggered false
        negatives. If instead the extrapolated mass exceeds the gas
        mass, it is similarly reverted to the 0th-order value (or, in a
        pathological case, hard-clamped to the gas mass).
  \item \textbf{\code{callloc=2}}, after the force loop
        (\code{chemistry\_end\_force}): a negative mass here means the
        limiters of Sections~\ref{sec:flux-limiter}--\ref{sec:anti-diffusion-guard}
        failed to prevent an overshoot. The particle's per-species
        \code{negativity\_counter} is incremented and a throttled warning
        is issued every $200$ consecutive occurrences
        (\code{GEAR\_FVPM\_DIFF\_NEGATIVITY\_COUNTER\_PRINT\_LIMIT}); if the
        magnitude of the negative mass ever exceeds the particle's gas
        mass, a hard \code{error()} is raised regardless of the
        counter, since that indicates a large, likely unrecoverable,
        conservation failure rather than round-off.
\end{itemize}

In addition to this three-site dispatcher:
\begin{itemize}
  \item \textbf{Final safety clamp} (Section~\ref{sec:state-reconstruction},
        step 6): an unconditional
        $U \gets \max(0,U)$ applied to both reconstructed states
        immediately before they are handed to the Riemann solver,
        independent of whether the softer \code{callloc=1} logic fired.
  \item \code{chemistry\_check\_unphysical\_metallicity}
        (\code{chemistry\_unphysical.h}) hard-clamps any
        reconstructed metallicity to $[0,1]$, used when the diffusion
        driver is $\Zfrac$.
  \item \code{chemistry\_check\_unphysical\_total\_metal\_mass}
        (\code{chemistry\_unphysical.h}) checks, after both the
        density loop and the force loop, that the \emph{sum} over all
        tracked species does not exceed the particle's gas mass; if it
        does, every species is rescaled by $1/(10\times m_{Z,\rm
        tot}/m)$ and a warning is issued (this is a diagnostic
        safety net, not expected to trigger in a correctly configured
        run).
  \item \code{chemistry\_get\_part\_corrected\_metal\_mass}
        (\code{chemistry\_getters.h}) is the \emph{only} interface
        exposed to other physics modules (feedback, cooling, star
        formation): it always clamps negative metal mass to zero and
        caps $m_Z \geq m$ down to $0.01\,m$. The diffusion solver is
        thus free to carry small transient negative or super-unity
        states internally (a deliberate design choice to avoid
        amplifying noise, see the comments in
        Section~\ref{sec:gradients}), while guaranteeing that no other
        module ever observes an unphysical metallicity.
\end{itemize}

\section{Timestep Stability Criteria}
\label{sec:timestep}

\subsection{GIZMO Particle Size}
The characteristic length entering every CFL estimate is the physical
GIZMO particle size,
\begin{equation}
  \dx = a \left(\frac{V_i}{\Omega_d}\right)^{1/d} ,
\end{equation}
$\Omega_d$ the volume of the $d$-dimensional unit sphere
(\code{chemistry\_get\_physical\_particle\_size}, \code{chemistry\_timesteps.h}
-- the same getter the hyperbolic solver's causal bound uses, so the CFL
length and the bound's length scale cannot drift apart); optionally, under
\code{GIZMO\_LANSON\_VILA\_PARTICLE\_SIZE}, the Lanson \& Vila (2008)
alternative $\dx = a/(2\,\overline{\delta x^{-1}})$ is used instead,
built from the same face-area accumulator as
Eq.~\eqref{eq:face_area}.

Their eq.~(58) defines $\dx_i = 1/(2\sum_l w_l\|A_{il}\|)$ with $w_l$ the
neighbour's \emph{volume}, evaluated here as $\sum_l\|A_{il}\|/V_l$; the
kernel value $w_jh_j^{-d}$ is dimensionally interchangeable with $1/V_l$
but smaller by a factor of order the neighbour count, and using it in that
place inflates $\dx_i$ by the same factor. The authors note below eq.~(58)
that $\sum_l w_l\|A_{il}\| = O(1/\epsilon_0)$, with $\epsilon_0$ the
smoothing length, so the two definitions describe the same resolution
scale and differ only by an $O(1)$ constant. Switching between them at a
fixed $\CCFL$ therefore buys nothing by itself: the constant is what a
Courant number absorbs, so the admissible $\CCFL$ has to be re-derived for
whichever length is in use, from Lanson \& Vila's own operator-based
stability condition (their eq.~62) rather than carried over.

\subsection{Mode-Dependent CFL Conditions}
\label{sec:cfl-modes}
The three diffusion modes do \emph{not} share a single timestep formula
(\code{chemistry\_diffusion\_timestep}, \code{chemistry\_timesteps.h}).
In every mode the Frobenius norm $\norm{\Kmat}_F$ is evaluated first; if
it vanishes (no diffusivity yet computed, e.g.\ at start-up) the
timestep constraint is skipped ($\dt\to\infty$).

\paragraph{Isotropic constant.}
\begin{equation}
  \label{eq:dt_const}
  \dt \le \CCFL\, \frac{\dx^2}{\norm{\Kmat}_F} = \CCFL\, \frac{\dx^2}{\kappaC} ,
\end{equation}
the textbook explicit-diffusion CFL bound (derived rigorously in
Appendix~\ref{app:von-neumann}), with no additional density-dependent
prefactor since $\driver=\rhoZ$ already carries the density.

\paragraph{Isotropic Smagorinsky and anisotropic gradient.}
Both go through the more elaborate form (Hopkins~\citeyear{Hopkins2017diffusion}, eq.~15-style)
\begin{equation}
  \label{eq:dt_aniso}
  \dt \le \CCFL\, \frac{E^2}{\norm{\Kmat}_F}\,\bar\rho ,
\end{equation}
where the extra factor $\bar\rho$ arises because here the conserved
state is $U=\rhoZ$ but the driver is $\driver=\Zfrac$, so
$\norm{U}/\norm{\driver} = \bar\rho$ (physical gas density,
\code{hydro\_get\_physical\_density}), and $E$ is a gradient-strength-aware
effective length,
\begin{equation}
  \label{eq:E-bracket}
  E =
  \begin{cases}
    \dx & \text{isotropic\_smagorinsky (see remark below)} \\[4pt]
    \dfrac{\norm{\driver}\,\dx}{\norm{\grad\driver}\,\dx + \norm{\driver}}
      & \text{anisotropic\_gradient,}
  \end{cases}
\end{equation}
with $\norm{\driver} = \big(\sum_{\rm species} \Zfrac^2\big)^{1/2}$ and
$\norm{\grad\driver}$ the Frobenius norm of the stacked
species-by-component metallicity gradient matrix, both evaluated
particle-locally. If $\norm{\driver}=0$ (no metal to diffuse) the
constraint is skipped.

For the anisotropic mode, $E$ interpolates between two regimes: in a
smooth region ($\norm{\grad\driver}\,\dx \ll \norm{\driver}$),
$E \to \dx$ and Eq.~\eqref{eq:dt_aniso} reduces to the same
$\dx^2/\norm{\Kmat}_F$ scaling as Eq.~\eqref{eq:dt_const}. In a region
with a locally steep metallicity gradient
($\norm{\grad\driver}\,\dx \gg \norm{\driver}$),
$E \to \norm{\driver}/\norm{\grad\driver}$, the gradient's own local
scale length, which can be much shorter than $\dx$; the timestep
constraint then becomes correspondingly \emph{tighter} than the naive
particle-size estimate. This is intentional and physically necessary --
diffusion locally erases steep features fastest exactly where they are
steep -- but the code comment at \code{chemistry\_timesteps.h}
notes explicitly that it is expensive in practice: strong feedback-driven
enrichment gradients can drive this term to very small values and
"slow down everything." For that reason the Smagorinsky mode
deliberately \emph{does not} use the bracket of
Eq.~\eqref{eq:E-bracket}, substituting the plain particle size $\dx$
instead (\code{chemistry\_timesteps.h}) -- a pragmatic,
non-rigorous compromise flagged as such directly in the source, adopted
because the physically correct bound was found to be prohibitively
restrictive for that closure in practice.

\subsection{Hyperbolic Relation}
\label{sec:hyperbolic-relation}
For completeness: under
\code{CHEMISTRY\_GEAR\_FVPM\_HYPERBOLIC\_DIFFUSION}, the same code base
instead advances a four-component relaxation system whose flux obeys a
Cattaneo-type equation with relaxation time $\tau$ and whose timestep
uses an ordinary hyperbolic CFL condition
$\dt = \CCFL\, \dx/v_{\max}$. As $\tau\to0$ the hyperbolic system's
flux equation collapses onto Eq.~\eqref{eq:ficks_law}, i.e.\ the
parabolic equation documented in this file is the relaxation limit of
that system; the two share the FVPM geometry, gradient estimator and
slope limiters described above, but use different Riemann solvers,
different conserved-state vectors, and different timestep criteria. The
hyperbolic solver's own physics (relaxation-time model, blending
factor between its HLL flux and the parabolic flux of
Eq.~\eqref{eq:flux_hll}) is documented separately and is not repeated
here.

\section{Boundary Conditions and Special Cases}

\subsection{No Explicit Boundary Treatment}
The module implements \emph{no} boundary-condition logic of its own --
no reflective, inflow/outflow, or fixed-concentration boundary is coded
anywhere in \code{GEAR\_FVPM\_DIFFUSION}. Periodicity, where enabled, is
resolved upstream by the SWIFT domain/tree infrastructure: the distance
vector $\mathbf{dx}_{ij}$ passed into every gradient and flux routine
documented above is already the minimum-image (periodicity-corrected)
separation by the time it reaches this module (see e.g.\ the comment at
\code{chemistry\_gradients.h}, "\emph{Do it this way in case dx
contains periodicity corrections already}"). For non-periodic domains,
no special treatment of domain edges is applied beyond the generic
neighbour-search behaviour of the host code.

\subsection{Isolated Particles}
The only edge case handled explicitly is a particle with zero
neighbours, for which \code{chemistry\_end\_density} is never called (so
the filtered fields would otherwise remain uninitialised).
\code{chemistry\_part\_has\_no\_neighbours} (\code{chemistry.h-...})
instead sets sane fallback values: the filtered density is set equal to
the particle's own local density and the filtered momentum to
$\bar\rho\,\mathbf v$, i.e.\ the particle is treated as its own,
unfiltered neighbourhood. Since $\Splus$ built from a single-particle
"gradient" is identically zero, such a particle contributes no
anisotropic diffusion until it acquires neighbours; the isotropic modes
are unaffected since $\kappaC$ does not depend on the shear tensor.

\section{Validation and Verification}

\subsection{Conservation Properties}
The pairwise sign convention of Section~\ref{sec:time-integration}
(flux subtracted from the source, added to the sink, both particles
updated from a single Riemann solve) guarantees exact discrete
conservation of total metal mass across every interacting pair,
independent of the flux-limiter stage of Section~\ref{sec:flux-limiter}
(the limiter rescales the flux by a common factor \emph{before} it is
split between the two particles, so the pairwise anti-symmetry is
preserved). Global conservation can therefore only be broken by the
\code{callloc=2} corrections of Section~\ref{sec:unphysical}, which are
intentionally non-conservative safety nets and are designed to trigger
rarely (tracked via the negativity counter for exactly this reason).

\subsection{Stability}
Section~\ref{sec:timestep}'s CFL bounds are necessary but, because of
the explicit minmod/artificial-diffusivity blending in
Eq.~\eqref{eq:flux_hll} and the nonlinear flux limiter of
Section~\ref{sec:flux-limiter}, not proven sufficient in closed form for
the full nonlinear scheme; Appendix~\ref{app:von-neumann} derives the
linear (constant-coefficient, unlimited) stability bound that motivates
Eq.~\eqref{eq:dt_const}.

\subsection{Suggested Test Problems}
\begin{itemize}
  \item \textbf{1D/3D diffusion of an isolated metal peak}
        (\code{isotropic\_constant} mode): compare the discrete solution
        against the analytic Green's function of the linear heat
        equation, Eq.~\eqref{eq:pde_const}, on a uniform particle
        distribution.
  \item \textbf{Two-peak merging test}: verify that two well-separated
        enriched regions diffuse independently at early times and merge
        smoothly, exercising the cell and face limiters of
        Section~\ref{sec:gradients} in a regime with two simultaneous
        local extrema.
  \item \textbf{Shear-flow anisotropic diffusion}: impose a prescribed
        background shear (e.g.\ a Couette or Kolmogorov flow) and verify
        that metal transport is preferentially aligned with the
        positive-eigenvalue directions of $\Splus$
        (Proposition~\ref{prop:psd}), not merely isotropic.
\end{itemize}

\section{Conclusion}

\subsection{Summary}
The parabolic branch of GEAR's FVPM chemistry module solves one of three
distinct diffusion PDEs (Eqs.~\eqref{eq:pde_const}--\eqref{eq:pde_aniso})
depending on the chosen closure, using a meshless least-squares gradient
estimator shared with the host MFM hydrodynamics solver, a two-stage
(cell- then face-wise) slope limiter tuned to avoid corrupting the
driver gradients that define the PDE itself, a Hopkins (2017)
minmod-blended HLL Riemann solver with an artificial-diffusivity cap and
an independent sign-consistency guard, and a physically-motivated,
smoothly-saturating flux limiter that enforces mass conservation bounds
without a hard discontinuity. Mode-dependent CFL conditions, several
layers of unphysical-state detection, and an unconditional final
positivity clamp round out the scheme's robustness.

\subsection{Limitations}
The explicit time integration ties the diffusion timestep to
Eq.~\eqref{eq:dt_const} or~\eqref{eq:dt_aniso}, which can become
severely restrictive in regions of high resolution or strong turbulence
(as the code's own comments acknowledge for the anisotropic bracket of
Eq.~\eqref{eq:E-bracket}). The Smagorinsky-mode timestep
(Section~\ref{sec:cfl-modes}) is a known, explicitly-flagged compromise
rather than a rigorous bound. The sign-consistency guard of
Section~\ref{sec:anti-diffusion-guard} can locally stall diffusion for a
sub-step whenever the HLL and direct flux estimates disagree, which is a
robustness trade-off rather than a formally analysed limiter property.

\subsection{Relation to the Wider Module}
This document deliberately excludes the hyperbolic (Cattaneo-relaxation)
formulation available under
\code{CHEMISTRY\_GEAR\_FVPM\_HYPERBOLIC\_DIFFUSION}
(Section~\ref{sec:hyperbolic-relation}); readers needing that solver's
four-component Riemann problem, relaxation-time model, and
parabolic/hyperbolic flux-blending factor should consult its own
documentation.

\appendix

\section{Mathematical Derivations}

\subsection{From Conservation and Fick's Law to the Parabolic Equation}
\label{app:derivation-pde}
Starting from the conservation law for a generic scalar density
$\rho Q$ advected with the flow and subject to a diffusive current
$\Fdiff$,
\begin{equation}
  \pdv{(\rho Q)}{t} + \div(\rho Q\,\mathbf v) = -\div\Fdiff ,
\end{equation}
the Lagrangian (particle-following) derivative satisfies
$\dv{}{t}(\rho Q) \equiv \pdv{(\rho Q)}{t} + \div(\rho Q\,\mathbf v) - \rho Q\,\div\mathbf v$
only up to the usual continuity-equation bookkeeping; in the FVPM
formalism this is sidestepped entirely because the conserved variable is
carried \emph{by} the moving, mass-conserving particle rather than
evolved on a fixed grid, so the advective flux never needs to be
discretised for a passive scalar: only the diffusive part remains to be
integrated,
\begin{equation}
  \pdv{(\rho Q)}{t}\bigg|_{\rm particle} = -\div\Fdiff .
\end{equation}
Substituting the tensorial Fick's law $\Fdiff = -\Kmat\cdot\grad q$
(Eq.~\eqref{eq:ficks_law}) gives the general parabolic equation
\begin{equation}
  \pdv{(\rho Q)}{t} = \div\big(\Kmat\cdot\grad q\big) ,
\end{equation}
which specialises to Eqs.~\eqref{eq:pde_const}--\eqref{eq:pde_aniso}
according to whether $q=\rho Q$ or $q=Q$ and whether $\Kmat$ is a scalar
multiple of the identity or the full anisotropic tensor $\kappaP\Splus$.

\subsection{Von Neumann Stability of the Explicit Parabolic Update}
\label{app:von-neumann}
Consider the 1D, constant-coefficient model problem
$\partial_t u = \kappaC\,\partial_x^2 u$ discretised with forward Euler
in time and central differences in space on a uniform grid of spacing
$\dx$ (the FVPM scheme reduces to this discretisation locally, in a
region of uniform particle spacing and constant $\Kmat=\kappaC\mathbf I$,
which is precisely the regime the code's
Eq.~\eqref{eq:dt_const} targets):
\begin{equation}
  \frac{u_i^{n+1}-u_i^n}{\dt} = \kappaC\,\frac{u_{i+1}^n - 2u_i^n + u_{i-1}^n}{\dx^2} .
\end{equation}
\begin{proposition}[Parabolic CFL bound]
The scheme above is Von Neumann stable if and only if
$\dt \le \dfrac{\dx^2}{2\kappaC}$.
\end{proposition}
\begin{proof}
Insert a Fourier mode $u_i^n = \hat u^n e^{\mathrm i k x_i}$, $x_i = i\,\dx$:
\begin{equation}
  \hat u^{n+1} = \hat u^n\Big[1 + \frac{\kappaC\dt}{\dx^2}\big(e^{\mathrm i k\dx} - 2 + e^{-\mathrm i k\dx}\big)\Big]
  = \hat u^n\Big[1 - \frac{4\kappaC\dt}{\dx^2}\sin^2\!\Big(\frac{k\dx}{2}\Big)\Big] ,
\end{equation}
using $e^{i\theta}+e^{-i\theta}-2 = -4\sin^2(\theta/2)$. The amplification
factor is
$G(k) = 1 - \dfrac{4\kappaC\dt}{\dx^2}\sin^2\!\big(\tfrac{k\dx}{2}\big) \in
\Big[1-\dfrac{4\kappaC\dt}{\dx^2},\ 1\Big]$
as $\sin^2(\cdot)$ ranges over $[0,1]$. Stability requires
$|G(k)|\le 1$ for every admissible $k$; since $G(k)\le1$ always, the
binding constraint is the lower bound,
\begin{equation}
  1 - \frac{4\kappaC\dt}{\dx^2} \ge -1
  \quad\Longleftrightarrow\quad
  \dt \le \frac{\dx^2}{2\kappaC} ,
\end{equation}
attained in the worst case $k\dx=\pi$ (the grid (particle)-scale
Nyquist mode). This is both necessary (violating it makes $|G(\pi/\dx)|>1$,
so that mode grows without bound under repeated application) and
sufficient (it makes $|G(k)|\le1$ for every $k$, so no mode grows).
\end{proof}
The code's constant $\CCFL$ (\code{chem\_data->C\_CFL\_chemistry}) absorbs
the factor of $1/2$ above together with an additional safety margin for
the (nonlinear, multi-dimensional, variable-coefficient, explicitly
minmod-limited) production scheme, for which the linear analysis above
is necessary but not sufficient. In the anisotropic case, replacing
$\kappaC \to \norm{\Kmat}_F$ in the bound is the standard generalisation
used when the exact eigenvalue decomposition of $\Kmat$ is not tracked
per-timestep for performance reasons: since
$\norm{\Kmat}_F \ge \lambda_{\max}(\Kmat)$ for any symmetric PSD matrix
(the Frobenius norm bounds the spectral radius from above), using
$\norm{\Kmat}_F$ in place of the true largest eigenvalue in
Eq.~\eqref{eq:dt_const}/\eqref{eq:dt_aniso} yields a \emph{more}
restrictive (i.e.\ still stable, possibly conservative) timestep bound.

\section{Nomenclature}

\subsection{Symbol to Code-Identifier Map}
\begin{center}
\begin{tabular}{@{}lll@{}}
\toprule
Symbol & Code identifier & Meaning \\
\midrule
$\rho$ & \code{hydro\_get\_physical\_density(p)} & Physical gas density \\
$\Zfrac$ & \code{chemistry\_get\_metal\_mass\_fraction} & Metal mass fraction \\
$\rhoZ$ & \code{chemistry\_get\_physical\_metal\_density} & Physical metal mass density \\
$m_Z$ & \code{p->chemistry\_data.metal\_mass[m]} & Conserved metal mass \\
$\driver$ & -- (selected by mode) & Diffusion driver: $\rhoZ$ or $\Zfrac$ \\
$\kappaC$ & \code{chem\_data->diffusion\_coefficient} & Global diffusion normalisation \\
$\kappaP$ & \code{p->chemistry\_data.kappa} & Per-particle diffusion coefficient \\
$\Kmat$ & \code{chemistry\_get\_physical\_matrix\_K} & Diffusion tensor \\
$\Kstar$ & \code{chemistry\_riemann\_compute\_K\_star} & Interface-averaged diffusion tensor \\
$\Shear$ & \code{chemistry\_get\_physical\_shear\_tensor} & Traceless velocity shear tensor \\
$\Splus$ & \code{chemistry\_regularize\_shear\_tensor} & PSD-regularised shear tensor \\
$\Fdiff$ & \code{chemistry\_get\_physical\_parabolic\_flux} & Diffusive flux \\
$\gammak$ & \code{kernel\_gamma} & Kernel compact-support factor \\
$\mathbf B_i$ & \code{p->geometry.matrix\_E} (already inverted) & Least-squares gradient matrix \\
$V_i$ & \code{p->geometry.volume} & Particle effective volume \\
$\mathbf A_{ij}$ & (local, \code{chemistry\_iact.h}) & Generalised interface area vector \\
$\xij$ & (local, \code{chemistry\_iact.h}) & Interface position relative to $i$ \\
$\psiHLL$ & \code{chem\_data->hll\_riemann\_solver\_psi} & HLL transport-term weight \\
$\epsHLL$ & \code{chem\_data->hll\_riemann\_solver\_epsilon} & Direct-flux sign-consistency tolerance \\
$\alpha$ & \code{chemistry\_riemann\_compute\_alpha} & Artificial diffusivity \\
$\CCFL$ & \code{chem\_data->C\_CFL\_chemistry} & Chemistry CFL coefficient \\
$\dx$ & (local, \code{chemistry\_timesteps.h}) & GIZMO physical particle size \\
\bottomrule
\end{tabular}
\end{center}

\subsection{Subscripts and Superscripts}
\begin{itemize}[noitemsep,topsep=0pt]
  \item $i,j$: particle indices (also used as tensor component indices where unambiguous from context)
  \item $L,R$: left/right Riemann states
  \item $n$: explicit timestep index
  \item $\star$: interface-averaged quantity
  \item $+$: positive-eigenvalue-projected (regularised) tensor
\end{itemize}

\subsection{Abbreviations}
\begin{itemize}[noitemsep,topsep=0pt]
  \item FVPM: Finite Volume Particle Method
  \item MFM: Meshless Finite Mass (GIZMO family of schemes)
  \item HLL: Harten--Lax--van Leer (approximate Riemann solver)
  \item CFL: Courant--Friedrichs--Lewy
  \item LES: Large-Eddy Simulation
  \item PSD: Positive Semi-Definite
\end{itemize}

\bibliography{references}

\end{document}
